\subsection{由正部分埃尔米特形式定义的自伴算子}

在本节中，E 表示一个复希尔伯特空间。

定义 8. --- 设 D 是 E 的向量子空间。E 上以 D 为定义域的部分埃尔米特形式是指一个对应 \(q = (\Gamma, \mathrm{E} \times \mathrm{E}, \mathbf{C})\) ，其定义域为 \(\mathrm{D} \times \mathrm{D}\) ，其图像 \(\Gamma\) 是函数性的，且从 \(\mathrm{D} \times \mathrm{D}\) 到 \(\mathbf{C}\) 中的、由 \(\Gamma\) 所定义的映射是埃尔米特形式。我们记 \(\operatorname{dom}(q) = \mathrm{D}\) 。

若 \(q\) 所定义的埃尔米特形式是正的，则称它为正部分形式。

设 u 是 E 上的自伴部分算子。对 E 的所有元素 x 与 y，我们用 \(\mu_{x,y}\) （相应地，\(\mu_{x}\) ）表示 \((x,y)\) （相应地，x）关于 u 的谱测度。

设 \((j, |u|)\) 是 \(u\) 的极分解（IV, p. 290 的定义 7）。用 D' 表示 \(|u|^{1/2}\) 的定义域。按定义，它是这样的 \(x \in E\) 所成的空间：函数 \(z \mapsto |z|\) 在 \(\mathrm{Sp}(u)\) 上关于测度 \(\mu_x\) 可积。它包含 \(\mathrm{dom}(u)\) 。

设 \((x,y)\in\mathrm{D}^{\prime}\times\mathrm{D}^{\prime}\) 。\(\mathrm{Sp}(u)\) 的恒等映射关于测度 \(\mu_{x,y}\) 可积（IV, p. 271 的命题 4, b)）；我们置

\[
q _ {u} (x, y) = \int_ {\operatorname{Sp} (u)} t d \mu_ {x, y} (t).
\]

若 \(y \in \operatorname{dom}(u)\) ，则由 IV, p. 271 的公式 (6)，我们有 \(q_u(x, y) = \langle x \mid u(y) \rangle\) 。映射 \(q_u\) 是 D' 上的埃尔米特形式。它定义了 E 上的一个部分埃尔米特形式，称为与 u 相伴的部分埃尔米特形式。若算子 u 是正的，则形式 \(q_u\) 是正部分埃尔米特形式。

定义 9. --- 设 \(q\) 是 E 上的正部分埃尔米特形式。我们用 \(\mathrm{E}_q\) 表示赋以下述纯量积的分离预希尔伯特空间 \(\operatorname{dom}(q)\)

\[
(x \mid y) _ {q} = \langle x \mid y \rangle + q (x, y).
\]

我们用 \(\|x\|_{q}\) 表示 \(x \in E_{q}\) 的范数。若 \(E_{q}\) 是希尔伯特空间，则称形式 q 是闭的。

命题 13. --- 设 u 是 E 上的正自伴部分算子。设 \(q_{u}\) 是与 u 相伴的正部分形式。

\begin{enumerate}
\def\labelenumi{\alph{enumi})}
\item
  \(q_{u}\) 的定义域是 \(u^{1/2}\) 的定义域，且对所有 \(\operatorname{dom}(u^{1/2})\) 中的 x 与 y，我们有 \(q_{u}(x,y)=\langle u^{1/2}(x)\mid u^{1/2}(y)\rangle\) ；
\item
  正部分形式 \(q_{u}\) 是闭的。u 的定义域是 \(q_{u}\) 的核。
\end{enumerate}

由于 \(u\) 是正的，我们有 \(|u| = u\) 。因此 \(\mathrm{dom}(|u|^{1/2})\) （即 \(q_u\) 的定义域）与 \(\mathrm{dom}(u^{1/2})\) 重合，并且我们有

\[
\begin{array}{c} q _ {u} (x, y) = \int_ {\operatorname{Sp} (u)} t d \mu_ {x, y} (t) = \int_ {\operatorname{Sp} (u)} t ^ {1 / 2} \cdot t ^ {1 / 2} d \mu_ {x, y} (t) \\ = \langle u ^ {1 / 2} (x) | u ^ {1 / 2} (y) \rangle \end{array}
\]

对 \(x\) 与 \(y\) （属于 \(\mathrm{dom}(q_u)\) ）成立（IV, p. 271 的命题 4, b)）。这就证明了断言 \(a\) )。

此外，预希尔伯特空间 \(E_{q_{u}}\) 这时与预希尔伯特空间 \(E_{u^{1/2}}\) 重合，后者相伴于 \(u^{1/2}\) （IV, p. 230 的定义 4）。由于 \(u^{1/2}\) 是闭部分算子，这个空间是希尔伯特空间（IV, p. 230 的命题 4），且 \(\text{dom}(u)\) 在 \(E_{u^{1/2}}\) 中稠密（依据 IV, p. 273 的推论 1 的断言 \(e\) )）。

设 \(q\) 是 \(E\) 上的部分埃尔米特形式，其定义域 \(D\) 在 \(E\) 中稠密。对 \(y \in D\) ，用 \(\lambda_y\) 表示线性形式 \(x \mapsto q(y, x)\) （定义在 \(D\) 上）。我们用 \(\widetilde{D}\) 表示这样的 \(y \in D\) 所成的集：线性形式 \(\lambda_y\) 在 \(D\) 上连续。

设 \(y \in \widetilde{D}\) 。由于 D 在 E 中稠密，存在 E 上唯一的连续线性形式，它延拓 \(\lambda_{y}\) 。我们仍用 \(\lambda_{y}\) 表示它。存在 E 中唯一的元素 \(u(y)\) ，使得 \(\lambda_{y}(x) = \langle u(y) | x \rangle\) 对所有 \(x \in E\) 成立（EVT, V, p. 15, 定理 3）。映射 \(y \mapsto u(y)\) 从 \(\widetilde{D}\) 到 E 中是线性的；我们称 E 上以 \(\widetilde{D}\) 为定义域的部分算子 u 为表示 q 的部分算子。因此我们有

\[
q (x, y) = \langle x \mid u (y) \rangle
\]

对 \(y \in \operatorname{dom}(u)\) 及 \(x \in D\) 成立。

注记. --- 设 \(q\) 是闭正部分形式，\(q'\) 是 E 上的连续正埃尔米特形式。在 \(\text{dom}(q)\) 上由 \((x, y) \mapsto q(x, y) + q'(x, y)\) 定义的正埃尔米特形式是以与 \(q\) 相同的定义域为定义域的闭正部分埃尔米特形式。我们把它记作 \(q + q'\) 。

依据 EVT, V, p. 16, 推论 2，存在唯一的线性映射 \(u' \in \mathcal{L}(\mathrm{E})\) ，使得 \(q'(x,y) = \langle x \mid u'(y) \rangle\) 对所有 \((x,y) \in \mathrm{E} \times \mathrm{E}\) 成立。自同态 \(u'\) 是正的，而表示闭正部分埃尔米特形式 \(q + q'\) 的部分算子是 \(u + u'\) 。

命题 14. --- 设 q 是 E 上稠定闭正部分形式，u 是表示 q 的部分算子。

\begin{enumerate}
\def\labelenumi{\alph{enumi})}
\item
  u 的定义域在希尔伯特空间 \(E_{q}\) 中稠密；
\item
  部分算子 u 是自伴且正的。
\end{enumerate}

由于 q 是闭的，定义 9 中的预希尔伯特空间 \(E_{q}\) 是希尔伯特空间。

我们来证明部分算子 \(u + 1_{E}\) （以 \(\operatorname{dom}(u)\) 为定义域）是双射的。由于

\[
\langle x \mid (u + 1 _ {\mathrm{E}}) (x) \rangle = q (x, x) + \| x \| ^ {2} \geqslant \| x \| ^ {2}
\]

对所有 \(x \in \operatorname{dom}(u)\) 成立，这个部分算子是单射的。

设 \(y \in E\) 。\(E_{q}\) 上由 \(x \mapsto \langle y | x \rangle\) 所定义的线性形式是连续的，因为 \(\|x\| \leqslant \|x\|_{q}\) 。因此存在 \(y' \in E_{q}\) 使得

\[
\langle y \mid x \rangle = (y ^ {\prime} \mid x) _ {q} = \langle y ^ {\prime} \mid x \rangle + q (y ^ {\prime}, x)
\]

对所有 \(x \in E_{q}\) 成立（EVT, V, p. 15, 定理 3）。按定义，这意味着 \(y'\) 属于部分算子 \(\widetilde{u}\) 的定义域，该算子表示以 dom(q) 为定义域的部分形式 \((x, y) \mapsto (x \mid y)_{q}\) ，且 \(\widetilde{u}(y') = y\) 。由于由上面的注记有 \(\widetilde{u} = u + 1_{E}\) ，我们由此推出部分算子 \(u + 1_{E}\) 也是满射的，从而是双射的。

我们来证明 \(u\) 的定义域在 \(\mathrm{E}_q\) 中稠密。设 \(y \in \mathrm{E}_q\) 与 \(\operatorname{dom}(u)\) 正交（在 \(\mathrm{E}_q\) 中）。存在 \(y' \in \operatorname{dom}(u)\) 使得 \(u(y') + y' = y\) 。这时我们有

\[
0 = (y \mid y ^ {\prime}) _ {q} = \langle y \mid y ^ {\prime} \rangle + q (y, y ^ {\prime}) = \langle y \mid y ^ {\prime} + u (y ^ {\prime}) \rangle = \| y \| ^ {2}
\]

因此 \(y = 0\) 。于是断言 \(a\) ) 得证。

由于 \(\mathrm{dom}(q)\) 在 E 中稠密，且 \(\| x\| \leqslant \| x\| _q\) 对所有 \(x\in \mathrm{dom}(q)\) 成立，断言 \(a\) ) 蕴含 \(\mathrm{dom}(u)\) 在 E 中稠密。

由于形式 \(q\) 是埃尔米特的（相应地，正的），对所有 \(x\) 与 \(y\) （属于 \(\operatorname{dom}(u)\) ），我们有

\[
\langle y \mid u (x) \rangle = q (y, x) = \overline {{q (x , y)}} = \overline {{\langle x \mid u (y) \rangle}} = \langle u (y) \mid x \rangle ,
\]

（相应地，\(\langle x \mid u(x) \rangle = q(x, x) \geqslant 0\) ），从而 u 是对称的（相应地，正的）。最后，由 IV, p. 240 的命题 10 的推论，部分算子 \(u + 1_{E}\) 是自伴的，u 亦然。

定理 3. --- 使 E 上的正自伴算子 u 对应于正部分形式 \(q_u\) 的映射，是 E 上正自伴部分算子所成的集与 E 上稠定闭正部分形式所成的集之间的一个双射。逆双射使正部分形式 q 对应于表示 q 的部分算子。

由命题 13, b) 与命题 14, b)，陈述中所描述的映射是妥善定义的。我们来证明它们互为逆双射。

设 \(u\) 是 E 上的正自伴部分算子，\(q\) 是与 \(u\) 相伴的正部分形式。用 \(v\) 表示表示 \(q\) 的正自伴算子。设 \(y \in \mathrm{dom}(u) \subset \mathrm{dom}(u^{1/2}) = \mathrm{dom}(q)\) 。对每个 \(x \in \mathrm{dom}(q) = \mathrm{dom}(u^{1/2})\) ，我们有

\[
q (y, x) = \langle u ^ {1 / 2} (y) | u ^ {1 / 2} (x) \rangle = \langle u (y) | x \rangle ,
\]

因此 y 属于 v 的定义域且满足 \(v(y) = u(y)\) 。于是部分算子 v 是 u 的一个延拓；由于 u 与 v 都是自伴的，它们相等。

反之，设 q 是稠定闭正部分形式，u 是表示 q 的正自伴部分算子。我们有 \(\operatorname{dom}(u) \subset \operatorname{dom}(u^{1/2})\) 。对 \(\operatorname{dom}(u)\) 中的 x 与 y，我们有

\[
q _ {u} (x, y) = \langle u ^ {1 / 2} (x) | u ^ {1 / 2} (y) \rangle = \langle x | u (y) \rangle = q (x, y)
\]

（依据命题 13, a)）。于是，希尔伯特空间 \(\mathrm{E}_q\) 与 \(\mathrm{E}_{q_u}\) 都包含 \(\operatorname{dom}(u)\) 作为稠密子空间（分别为命题 14, a) 与命题 13, b)），且 \(\|x\|_q = \|x\|_{q_u}\) 对所有 \(x \in \operatorname{dom}(u)\) 成立。由此得出 \(\mathrm{E}_q = \mathrm{E}_{q_u}\) 且 \(q = q_u\) （引理 8）。

推论. --- 设 \(q\) 是 \(E\) 上的闭正部分形式，而 \(u\) 是表示 \(q\) 的正自伴算子。\(q\) 的定义域等于 \((1_{E} + u)^{1/2}\) 的定义域，且我们有

\[
\| x \| _ {q} = \| (1 _ {\mathrm{E}} + u) ^ {1 / 2} x \|
\]

对所有 \(x \in \operatorname{dom}(q)\) 成立。

\(q\) 的定义域等于 \(u^{1/2}\) 的定义域（命题 13, \(a\) )），后者与 \((1_{\mathrm{E}} + u)^{1/2}\) 的定义域重合（参见 IV, p. 271 的命题 3）。对每个 \(x \in \operatorname{dom}(u) \subset \operatorname{dom}(u^{1/2})\) ，我们有

\[
\| (1 _ {\mathrm{E}} + u) ^ {1 / 2} x \| ^ {2} = \langle x | (1 _ {\mathrm{E}} + u) (x) \rangle = \| x \| ^ {2} + \langle x | u (x) \rangle = \| x \| _ {q} ^ {2}.
\]

由于 \(u\) 的定义域在希尔伯特空间 \(\mathrm{E}_q\) 中稠密（命题 14, a)），这个公式由连续性延拓到所有 \(x \in \mathrm{dom}(u^{1/2})\) 。

例. --- 设 u 是 E 上的正部分算子，它未必是闭的。用下式定义以 dom(u) 为定义域的正部分形式 q

\[
q (x, y) = \langle x \mid u (y) \rangle
\]

对 \(x\) 与 \(y\) （属于 \(\mathrm{dom}(u)\) ）成立。设 \(\mathrm{E}_q\) 是 IV, p. 292 的定义 9 中的分离预希尔伯特空间。设 \(\widetilde{\mathrm{E}}_q\) 是 \(\mathrm{E}_q\) 的完备化希尔伯特空间，其纯量积仍记作 \((x,y) \mapsto (x \mid y)_q\) 。由于 \(\iota\) （\(\mathrm{E}_q\) 到 E 中的典范包含）是连续的，它具有唯一的连续延拓，记作 \(\widetilde{\iota}\) （从 \(\widetilde{\mathrm{E}}_q\) 到 E 中）。埃尔米特形式 \(q\) 在 \(\mathrm{E}_q\) 上是连续的，因而也延拓为唯一的连续正埃尔米特形式 \(\widetilde{q}\) ，定义在 \(\widetilde{\mathrm{E}}_q\) 上。

我们来证明线性映射 \(\widetilde{\iota}\) 是单射的。设 \(x\in\mathrm{Ker}(\widetilde{\iota})\) 。考虑序列 \((x_{n})_{n\in\mathbf{N}}\) （属于 \(\mathrm{E}_{q}\) ），它收敛于 \(x\) ，在 \(\widetilde{\mathrm{E}}_{q}\) 中。这时我们有

\[
\lim _ {n \to + \infty} \iota (x _ {n}) = \widetilde {\iota} (x) = 0,
\]

因此序列 \((x_{n})_{n\in\mathbf{N}}\) 在 E 中收敛于 0。设 \(y\in E_{q}\) 。由于 \(\widetilde{q}\) 在 \(\widetilde{E}_{q}\) 上连续，由此可得

\[
\begin{array}{r l} (x \mid y) _ {q} = \lim _ {n \to + \infty} (x _ {n} \mid y) _ {q} = & \lim _ {n \to + \infty} \Big (\langle x _ {n} \mid y \rangle + q (x _ {n}, y) \Big) \\ & = \lim _ {n \to + \infty} \Big (\langle x _ {n} \mid y \rangle + \langle x _ {n} \mid u (y) \rangle \Big) = 0. \end{array}
\]

由于空间 \(E_{q}\) 在 \(\widetilde{E}_{q}\) 中稠密，我们推出 x = 0，此即所要证者。

把 \(\widetilde{E}_{q}\) 与它在 \(\widetilde{\iota}\) 下的像等同于 E 的子空间，我们便把 \(\widetilde{q}\) 解释为延拓 q 的稠定闭正部分形式。与 \(\widetilde{q}\) 相伴的正自伴算子是 u 的一个自伴延拓。它称为 u 的弗里德里希斯扩张。

注. --- 设 \(c \in \mathbf R_{+}\) 。设 q 是稠定闭部分埃尔米特形式，使得对所有属于 q 的定义域的 x 有 \(q(x, x) \geqslant -c \|x\|^2\) 。这意味着由 \(\widetilde{q}(x, y) = q(x, y) + c \langle x \mid y \rangle\) 定义的以 dom(q) 为定义域的闭部分埃尔米特形式是正部分形式。于是由定理 3，这样的形式对应于 E 上使 \(u + c1_{E}\) 为正的自伴部分算子 u，亦即使得 u 以 -c 为下界（IV, p. 237 的定义 7）者。

反之，设 u 是 E 上的对称部分算子，使得

\[
\langle x \mid u (x) \rangle \geqslant - c \| x \| ^ {2}\tag{17}
\]

对所有 \(x \in \operatorname{dom}(u)\) 成立。这时 \(v = u + c1_{E}\) 是 E 上的正部分算子。u 的弗里德里希斯扩张是自伴算子 \(\widetilde{v} - c1_{E}\) ，其中 \(\widetilde{v}\) 是 v 的弗里德里希斯扩张；它是 u 的一个自伴延拓，不依赖于满足 (17) 的实数 c 的选取。
% semantic-fragments: legacy-input-seam
\relax{}
